% ID: SEL-D0072-OA-E03
% ESTADO: validada. S072.
\textbf{Parte I. a)} Por la definición de probabilidad condicionada, $0{,}60=0{,}30/P(B)$, de donde $P(B)=\boxed{0{,}50}$. Por inclusión y exclusión, $1=P(A)+0{,}50-0{,}30$, luego $P(A)=\boxed{0{,}80}$.

\textbf{b)} $P(A)P(B)=0{,}40$ difiere de $P(A\cap B)=0{,}30$. Los sucesos son dependientes.

\textbf{Parte II. a)} El centro del intervalo normal es la media muestral:
\[
\bar x=(509{,}41+539{,}79)/2=\boxed{524{,}60\text{ euros}}.
\]
\textbf{b)} El margen impreso es $(539{,}79-509{,}41)/2=15{,}19$ euros. Al $97\%$ usamos $z_{0{,}985}\approx2{,}17009$. Al despejar de $E=z_{0{,}985}\,84/\sqrt n$ resulta $n\approx144{,}01$. El tamaño recuperado es $\boxed{144}$: con ese entero, el margen es $2{,}17009\cdot7\approx15{,}1906$ y reproduce ambos extremos a centésimas.
